% =====================================================================
\chapter{Discusión y conclusiones}
\label{cap:conclusiones}

% ---------------------------------------------------------------------
\section{La hipótesis}

La hipótesis sostenía que es posible construir una metodología de ciencia de datos que mida el
sentimiento de las noticias económicas argentinas y lo relacione con el comportamiento de las
acciones. \textbf{Se acepta}: se cumplen las dos partes del criterio fijado de antemano.

\begin{itemize}
  \item \textbf{Medir.} FinBETO clasifica oraciones de noticias que nunca vio con un F1 macro de
        \dato{fb.f1_macro_media@3}, \dato{fbcmp.completo.ventaja_puntos} puntos por encima de un
        modelo de sentimiento en español ya existente, que confunde casi todo con neutro.
  \item \textbf{Relacionar.} En ruedas que ningún modelo vio, el mercado mejora la predicción del
        sentimiento de las noticias de la rueda siguiente ($\Delta\text{AUC} =
        \dato{ps.d.logistica.delta_auc@s3}$). En la dirección contraria no hay nada: el
        sentimiento no mejora la predicción del precio ($\Delta\text{AUC} =
        \dato{np.d.logistica.delta_auc@s3}$), en ninguna de las nueve variantes de robustez.
\end{itemize}

La respuesta a la pregunta del título es, entonces, que en el período estudiado \textbf{los medios
económicos argentinos siguieron al mercado}: contaron lo que el mercado ya había hecho, y no lo
anticiparon. Coincide con lo que se esperaba, y con la parte de la evidencia internacional que
encuentra que el sentimiento de los medios refleja los precios \citep{tetlock2007,verma2026}.

% ---------------------------------------------------------------------
\section{Los objetivos}

\begin{table}[H]
\centering\small
\begin{tabularx}{\linewidth}{@{}C{0.9cm}YC{1.8cm}@{}}
\toprule
\textbf{Obj.} & \textbf{Resultado obtenido} & \textbf{¿Se cumplió?} \\
\midrule
1 & Corpus de \dato{corpus.notas} notas de \dato{corpus.portales} portales, auditado, con \dato{regla17.notas_en_ventana_union} notas asignadas a las \dato{dataset.ruedas} ruedas de la ventana & Sí \\
2 & FinBETO, de tres clases: F1 macro de \dato{fb.f1_macro_media@3} en noticias no vistas, contra \dato{fbcmp.completo.referencia_f1_macro@3} del modelo de referencia & Sí \\
3 & Sentimiento inesperado, sentimiento propio de \dato{dataset.empresas} empresas, retorno, volumen anormal y VIX argentino con su escala, en un dataset de \dato{dataset.filas_panel} filas & Sí \\
4 & Modelo predictivo evaluado en \dato{np.ruedas_test} ruedas no vistas: el mercado anticipa al sentimiento; el sentimiento no anticipa al mercado & Sí \\
\bottomrule
\end{tabularx}
\caption{Cumplimiento de los objetivos específicos. Fuente: elaboración propia.}
\label{tab:cumplimiento}
\end{table}

% ---------------------------------------------------------------------
\section{Qué quiere decir el resultado}

\textbf{Por qué el sentimiento no anticipa el precio.} Es lo que predice la hipótesis de mercados
eficientes \citep{fama1970}: si una noticia trae información, el precio la incorpora el mismo día,
antes de que sirva para predecir el día siguiente. Además, buena parte de lo que publican los
medios económicos es la crónica del propio mercado.

\textbf{Por qué el precio sí anticipa el sentimiento.} La variable que más pesa en el modelo es el
retorno del Merval del mismo día. Por la regla de las 17:00, las notas publicadas después del cierre
cuentan para la rueda siguiente, y son justamente las que cuentan cómo cerró el mercado: si el
Merval subió, se escribe de forma más positiva. La relación es el reflejo del mercado en la
prensa, no un efecto de la prensa sobre el mercado.

\textbf{Para qué sirve.} Para quien sigue el mercado argentino, leer el sentimiento de los medios
económicos como una señal de lo que va a pasar es leer un reflejo de lo que ya pasó. Para la
ciencia de datos, el trabajo deja un pipeline reproducible y un modelo de sentimiento para la
prensa económica en español rioplatense, que no existían.

% ---------------------------------------------------------------------
\section{Desaciertos y límites}
\label{sec:limites}

\textbf{Lo que salió mal en el camino y se corrigió.}
\begin{itemize}
  \item \textbf{La primera evaluación de FinBETO estaba inflada.} Las oraciones se habían dividido
        por oración y no por noticia, y había textos repetidos entre entrenamiento y prueba. Al
        dividir por noticia, el F1 bajó a su valor honesto. Es el error de fuga de información más
        común con texto \citep{kapoor2023leakage}.
  \item \textbf{La clase MIXTO no se usaba.} Se había creado una cuarta clase que el sentimiento
        neto trataba igual que la neutra. Se unió con NEU (sección~\ref{sec:tresclases}).
  \item \textbf{La cantidad de notas no era confiable.} El plan original definía el shock
        informativo como un aumento de la cantidad de noticias. Con un corpus recolectado de un
        archivo web, esa cantidad depende de qué se capturó, y se la dejó afuera de los modelos.
  \item \textbf{Demasiadas pruebas.} La versión anterior contrastó la hipótesis con cientos de
        pruebas de Granger y correcciones por pruebas múltiples. Se reemplazaron por una sola
        comparación predictiva, fijada de antemano, y Granger quedó como verificación.
\end{itemize}

\textbf{Los límites del resultado.}
\begin{itemize}
  \item \textbf{El test es corto y calmo.} Tiene \dato{np.ruedas_test} ruedas, y todas cayeron en
        el régimen Calmo: el VIX promedió \dato{vix.vix_media_test}\% en el test, contra
        \dato{vix.vix_media_entrenamiento}\% en el entrenamiento. El resultado vale para un mercado
        tranquilo; \citet{garcia2013} encontró que el sentimiento anticipa los retornos sobre todo
        en las recesiones, un período que esta ventana no tiene.
  \item \textbf{El efecto se achica con el sentimiento ajustado por portal.} Si a cada nota se le
        resta el promedio de su portal, el $\Delta$AUC de precio \flecha{} sentimiento baja a
        \dato{rob.sentimiento_por_portal.ps.logistica.delta_auc@s3} y su intervalo incluye el 0.
        Puede indicar que parte del efecto viene de \textit{qué} portales publican en cada rueda,
        por ejemplo, de los que hacen la crónica del cierre. Es una hipótesis que no se verificó.
  \item \textbf{El tamaño depende de la ventana.} Con 10 o con 60 ruedas el $\Delta$AUC sigue
        siendo positivo, pero su intervalo incluye el 0.
  \item \textbf{Las etiquetas no tienen revisión humana.} El F1 de FinBETO mide su acuerdo con Gemma,
        no con un criterio humano.
  \item \textbf{El corpus.} El \dato{regla17.pct_hora_sospechosa}\% de las notas tiene una hora
        por defecto, que puede mandarla a la rueda equivocada; abril de 2026 no tiene notas de
        iProfesional (sacarlo no cambia el resultado); y el sentimiento propio solo se pudo medir
        en \dato{dataset.empresas} empresas.
  \item \textbf{Los datos.} Un año de ruedas, datos diarios, que no ordenan lo que pasa dentro de
        una rueda, y precios en pesos nominales.
\end{itemize}

% ---------------------------------------------------------------------
\section{Líneas futuras}

Etiquetar a mano una muestra de oraciones para medir el acuerdo de Gemma con personas; recuperar el
matiz de las oraciones mixtas como una dimensión aparte; analizar el resultado portal por portal,
para explicar por qué se achica con el ajuste por portal; extender la ventana a varios años, con
períodos de crisis, para probar si el sentimiento anticipa al mercado cuando está tenso, como
encontró \citet{garcia2013}; repetir el modelo con precios intradiarios y en dólares; e incorporar
otras fuentes, como redes sociales.
